% =========================================
% CHAPTER 5: IMPLEMENTATION
% =========================================
\chapter{Implementation Details}
\label{ch:implementation}

\section{Overview}

This chapter describes how the framework is built. The design chapter (Chapter~\ref{ch:framework}) explained what we want to measure and why. This chapter explains what we actually wrote and what trade-offs we made.

Three things had to work together:

\begin{enumerate}
    \item \textbf{Faithful architecture.} The ZK contract must require a proof at submission time and enforce a post-submission confirmation delay. The Optimistic contract must accept submissions without a proof and enforce a challenge window before finality.

    \item \textbf{Shared Layer~1.} Both sequencers operate against the same Hardhat node simultaneously. This eliminates timing and block-production differences as confounds.

    \item \textbf{Instrumented measurement.} Every state transition records a timestamp and outcome so we can analyze the data later.
\end{enumerate}

The implementation has three layers: Layer~1 smart contracts (finality rules), Layer~2 sequencer processes (batching and submission), and a metrics stack (Prometheus and Grafana). Where the framework simplifies production behavior, we explain why and what that means for the results.

\section{Layer~1: Smart Contract Design}
\label{sec:l1-design}

Both contracts follow patterns established in production systems. The \texttt{ZKVerifier} draws on the Groth16 verification structure from Scroll (Chapter~\ref{ch:relatedwork}, Section~\ref{sec:zk-production}). The \texttt{OptimisticVerifier} draws on the multi-state finality model from Optimism Bedrock (Section 3.4 in the related work chapter).

\subsection{Core Design Difference: Commitment Structures}

The essential architectural difference between the two approaches manifests in the data each contract stores. For ZK, a commitment is final once its proof has been verified and a confirmation delay has elapsed. For Optimistic, a commitment enters a challenge period and becomes final only after that period expires without a successful dispute. Table~\ref{tab:contract-diffs} summarizes the data structures that implement this difference.

\subsection{ZKVerifier Contract}

The \texttt{ZKVerifier} contract is the on-chain anchor for the ZK~Rollup path. It accepts proofs submitted by the sequencer and enforces a finality delay before commitments are irreversible. The contract implements two core functions:

\begin{enumerate}
    \item \texttt{submit(bytes32 stateRoot, bytes calldata proof)} — Records a state commitment and its accompanying proof. In this framework, proof verification is mocked: the function checks that a proof is provided but does not execute the full BN254 pairing computation that a production Groth16 verifier would perform. (Production Scroll~\cite{scroll2023} implements the pairing check; we simplify to preserve experimental iteration speed.)

    \item \texttt{isFinalized(uint256 id)} — Returns true if the commitment has been recorded for at least five additional Layer~1 blocks. This simulates the confirmation delay in production systems before a commitment is cryptographically irreversible.
\end{enumerate}

The rationale for each design choice is as follows:

\textbf{Proof verification simplification}: Full BN254 pairing verification requires substantial gas (200K-500K per proof~\cite{groth2016size}). For this framework, we preserve the architectural semantics (proof is required at submission, no challenge window needed) without incurring the full computational cost. The sequencer still generates a proof; the contract still accepts and records it; the only simplification is the verification check itself.

\textbf{Five-block finality}: Production Ethereum-based rollups typically wait for 12-15 block confirmations before treating a commitment as final. We use 5 blocks in this framework to keep experimental runs tractable while preserving the order of magnitude of the delay. Any Layer~1 reorganization shorter than 5 blocks will not affect finality in either system.

\textbf{State structure}: The contract stores a \texttt{Commitment} struct per batch, containing the state root, submission block number, and proof hash. Events are emitted on every submission, allowing the sequencer and dashboard to track finality propagation without polling.

The full contract implementation is available in the framework repository at \texttt{l1/contracts/ZKVerifier.sol}.

\subsection{OptimisticVerifier Contract}

The \texttt{OptimisticVerifier} contract implements the fraud-proof challenge model that defines Optimistic~Rollups. It is drawn from patterns in Optimism Bedrock~\cite{optimism2023bedrock} and Arbitrum Nitro~\cite{nitro2023}. The contract implements four core functions:

\begin{enumerate}
    \item \texttt{submit(bytes32 stateRoot)} — Records a state commitment without a proof and opens a challenge window. The window deadline is computed as current block number plus \texttt{CHALLENGE\_PERIOD} (10 blocks in this framework).

    \item \texttt{challenge(uint256 id)} — Allows any observer to dispute a commitment within the challenge window. The contract records that a challenge has been filed; in this framework, the challenge permanently invalidates the commitment. (Production systems implement an interactive dispute game to determine if the challenge is valid; we simplify to a binary outcome.)

    \item \texttt{finalize(uint256 id)} — Makes explicit on-chain that a commitment has survived its challenge period without dispute. After calling this function, the commitment is marked \texttt{finalized}.

    \item \texttt{isFinalized(uint256 id)} — Returns true if the commitment is explicitly finalized OR if the challenge period has expired without a dispute. This dual path ensures that finality is recognized whether or not anyone calls the explicit \texttt{finalize()} function.
\end{enumerate}

\textbf{Challenge period duration}: We set the challenge period to 10 blocks (~2 minutes on a local Hardhat node). Production systems use 7 days~\cite{arbitrum-original}; the compressed window is necessary for experimental tractability while preserving the architectural semantics. Any test scenario can be executed and analyzed in seconds rather than days.

\textbf{Dispute resolution}: In production Arbitrum and Optimism, a valid fraud proof requires an interactive bisection game followed by on-chain re-execution of a single VM step. Here, we record a dispute but do not re-execute; the challenge simply invalidates the root. This simplification removes the largest source of complexity while preserving the core property: incorrect state roots can be challenged and prevented from reaching finality.

\textbf{Finality model}: The four-state model (\texttt{PENDING}, \texttt{SOFT\_FINAL}, \texttt{HARD\_FINAL}, \texttt{CHALLENGED}) inherited from Optimism Bedrock allows the dashboard and test suite to distinguish finality stages. This is essential for the comparative analysis in Chapter~\ref{ch:evaluation}: users and applications may behave differently if a commitment is soft-final (deadline passed, no challenge yet) versus hard-final (explicit on-chain finalization).

The full contract implementation is available in the framework repository at \texttt{l1/contracts/OptimisticVerifier.sol}.

\section{Layer~2: Sequencer Processes}
\label{sec:l2-design}

Both sequencers implement the same high-level logic: maintain a pending transaction pool, periodically create batches, compute commitments, and submit to Layer~1. The critical difference is in the commitment computation: ZK sequencers generate and submit a proof alongside the state root, while Optimistic sequencers submit only the state root. This difference drives all downstream architectural choices.

\subsection{Architectural Difference: Proof vs. No Proof}

The sequencers are implemented as independent Python processes using Flask~\cite{flask2024} for HTTP and web3.py~\cite{web3py2024} for Layer~1 interaction. Table~\ref{tab:sequencer-comparison} summarizes how the two differ in the commitment step.

\begin{table}[h]
    \centering
    \begin{tabular}{l|c|c}
        \hline
        & \textbf{ZK Sequencer} & \textbf{Optimistic Sequencer} \\
        \hline
        Commitment step & Two SHA-256 hashes & One SHA-256 hash \\
        \hline
        Batch threshold & 100 transactions & 200 transactions \\
        \hline
        Polling interval & 5 seconds & 8 seconds \\
        \hline
        L1 gas per submission & $\sim$200,000 (Hardhat only) & $\sim$150,000 (Hardhat only) \\
        \hline
        Contract loading & Per \texttt{\_batch()} call (lazy) & Once in \texttt{\_\_init\_\_} \\
        \hline
        L1 call & \texttt{ZKVerifier.submit(root, proof)} & \texttt{OptimisticVerifier.submit(root)} \\
        \hline
    \end{tabular}
    \caption{Architectural differences between sequencer types. Gas values are from a local Hardhat node and are not representative of Ethereum mainnet costs.}
    \label{tab:sequencer-comparison}
\end{table}

\textbf{Rationale for these differences}: The Optimistic model does not require proof generation at submission time, so the sequencer can afford larger batches and longer polling intervals. Batches of 200 transactions amortize the fixed gas cost of Layer~1 submission more effectively than 100-transaction batches. The longer polling interval (8 seconds vs. 5 seconds) reflects lower latency pressure: the Optimistic model's challenge window determines finality time, not submission latency.

Conversely, the ZK sequencer's proof generation imposes a computational bottleneck at submission. To keep finality latency low and allow users to see their transactions finalized quickly, the ZK sequencer uses smaller batches and shorter polling intervals. The choice of 100 transactions per batch and 5-second polling is calibrated to keep batch assembly time below 1 second on typical hardware, enabling finality within 5-10 seconds.

The gas allocations (200K vs. 150K) reflect the difference in Layer~1 contract work. The \texttt{ZKVerifier.submit()} function stores both a state root and a proof hash, while \texttt{OptimisticVerifier.submit()} stores only the root and opens a challenge period. Less on-chain work means lower gas cost.

\subsection{Sequencer Implementation Structure}

Both sequencers follow the same four-step structure:

\begin{enumerate}
    \item \texttt{add\_tx(tx)}: Accept a transaction via HTTP POST and append it to the pending pool.
    \item \texttt{periodic\_batch()}: Background thread that wakes up every N seconds; if the pool has reached the batch threshold, call \texttt{\_batch()}.
    \item \texttt{\_batch()}: Compute the state root (and proof, if ZK), then submit to Layer~1. The two sequencers differ here in two ways:
        \begin{itemize}
            \item The \textbf{ZK sequencer} reads \texttt{/app/contracts.json} on \emph{every} \texttt{\_batch()} call to obtain the verifier address. There is no contract handle stored in \texttt{\_\_init\_\_}. If the file is absent (for example, deployment has not completed), the exception is caught and the batch is dropped silently.
            \item The \textbf{Optimistic sequencer} calls \texttt{\_init\_contract()} in \texttt{\_\_init\_\_} and stores the contract handle as \texttt{self.contract}. If that initial call failed because the L1 node was not ready, \texttt{\_batch()} retries \texttt{\_init\_contract()} before proceeding.
        \end{itemize}
        This difference has a practical effect: the ZK sequencer tolerates a slow deploy by simply skipping batches until the file appears; the Optimistic sequencer needs a successful L1 connection at startup.
    \item \texttt{calculate\_tps()}: Background thread that updates the throughput gauge periodically.
\end{enumerate}

The Flask HTTP API exposes three endpoints:
\begin{itemize}
    \item \texttt{POST /tx} — Accept a transaction (returns 202 Accepted immediately)
    \item \texttt{GET /metrics} — Return current transaction/batch counts and throughput
    \item \texttt{GET /health} — Return liveness status (Optimistic sequencer only)
\end{itemize}

Both sequencers export Prometheus metrics on dedicated ports (9100 for ZK, 9101 for Optimistic), enabling the monitoring stack to observe finality latency, proof generation time, and other real-time data.

\textbf{Implementation simplification}: In production systems, the sequencer also performs state execution (applying transactions to a state tree, updating storage, etc.). This framework omits that step and uses synthetic state roots (SHA-256 hashes of batch data). This simplification allows the framework to run on modest hardware while preserving the key architectural difference: ZK requires proof generation at every submission, while Optimistic does not. The comparative analysis in Chapter~\ref{ch:evaluation} remains valid because it measures the differences that matter—finality latency, throughput, and proof generation time—regardless of whether state execution is simplified.

Both sequencer implementations are available in the framework repository: \texttt{l2-zk/sequencer.py} and \texttt{l2-optimistic/sequencer.py}.

\section{Web Dashboard and Metrics Infrastructure}
\label{sec:metrics-design}

The framework includes a Flask~\cite{flask2024}-based web dashboard that aggregates real-time metrics from both sequencers and exposes HTTP endpoints for test automation. The dashboard runs on port~3000 and provides:

\begin{itemize}
    \item \textbf{Real-time status}: Current transaction count, batch count, and throughput for each sequencer.
    \item \textbf{Comparative visualization}: Side-by-side graphs of finality latency and throughput over time, enabling visual identification of architectural differences.
    \item \textbf{Test control API}: Endpoints to submit test transactions, query sequencer health, and trigger adversarial scenarios (crash simulation, delay injection).
\end{itemize}

\begin{figure}[h]
    \centering
    \includegraphics[width=0.9\textwidth]{dashboards/flask-dashboard-overview.png}
    \caption{Flask dashboard showing real-time metrics from both sequencers (transaction counts,
    current TPS, batch totals, and throughput over time). \emph{[Screenshot to be captured
    per \texttt{ths/images/README.md} Step~1: run stack, trigger Normal~150 scenario on each
    sequencer, screenshot \texttt{http://localhost:3000}. Save as
    \texttt{ths/images/dashboards/flask-dashboard-overview.png}.]}
    }
    \label{fig:flask-dashboard}
\end{figure}

Underlying the dashboard is a Prometheus~\cite{prometheus2024} metrics collection layer. Each sequencer exports metrics to a dedicated HTTP endpoint (ports 9100 and 9101) that Prometheus scrapes at 15-second intervals. Table~\ref{tab:metrics-exported} lists the metrics collected.

\begin{table}[h]
    \centering
    \begin{tabular}{l|l|p{3cm}}
        \hline
        \textbf{Metric} & \textbf{Type} & \textbf{Interpretation} \\
        \hline
        l2\_transactions\_total & Counter & Cumulative count of accepted L2 transactions \\
        \hline
        l2\_batches\_total & Counter & Cumulative count of batches submitted to L1 \\
        \hline
        l2\_tps & Gauge & Rolling estimate of throughput (tx/sec) \\
        \hline
        l2\_finality\_time\_seconds & Histogram & Latency from L1 submission to finality, with percentiles (P50, P95, P99) \\
        \hline
        l2\_proof\_generation\_seconds & Histogram & Time to generate mock ZK proof (ZK sequencer only) \\
        \hline
        l2\_pending\_transactions & Gauge & Current size of pending transaction pool \\
        \hline
    \end{tabular}
    \caption{Prometheus metrics collected from sequencers}
    \label{tab:metrics-exported}
\end{table}

\textbf{Why this instrumentation matters}: The finality latency histogram answers RQ1—we can see exactly how long it takes for transactions to become irreversible, not just an average. The proof generation histogram isolates the computational cost of the ZK path. By comparing these histograms, we can tell whether ZK is faster because of architectural differences, or because our implementation is optimized differently. The pending transaction gauge shows how batching parameters affect how long users wait before their transaction gets included in a batch.

Grafana~\cite{grafana2024} connects to the Prometheus datasource and renders pre-configured dashboards automatically when the Docker Compose~\cite{compose2024} stack starts. This ensures that anyone deploying the framework sees the same visualizations without manual configuration.

\begin{figure}[h]
    \centering
    \includegraphics[width=0.9\textwidth]{dashboards/grafana-comparison.png}
    \caption{Grafana side-by-side comparison dashboard showing finality time, throughput, batch
    efficiency, and proof generation time (ZK only).
    \emph{[Screenshot to be captured per \texttt{ths/images/README.md} Step~2: after running
    High Load scenario, navigate to \texttt{http://localhost:3001}, open
    ``ZK / Optimistic Rollup L2 Metrics Comparison'', set time range to Last~5~minutes.
    Save as \texttt{ths/images/dashboards/grafana-comparison.png}.]}
    }
    \label{fig:grafana-comparison}
\end{figure}

\section{Summary: Implementation and Big Picture}

This chapter described how we translated architectural concepts into working code. The core principle is: keep the architecture faithful, simplify the expensive parts, and measure the differences that matter.

The Layer~1 contracts preserve the essential security models (ZK requires proof at submission; Optimistic requires a challenge window), but they don't do expensive cryptographic operations. The Layer~2 sequencers implement the core batching and submission logic that distinguishes ZK from Optimistic, with tuned parameters that reflect each model's economics. All parameters are configurable via environment variables, so researchers can change batch sizes, polling intervals, or gas allocations and immediately see the impact.

The metrics infrastructure captures what we need to answer the research questions:

\begin{itemize}
    \item RQ1 (Finality latency) — measured via the finality time histogram
    \item RQ2 (Recovery cost) — measured via adversarial test scenarios that trigger crashes
    \item RQ3 (Comparative evaluation) — enabled by running both sequencers against the same Layer~1 node
    \item RQ4 (Reproducibility) — achieved via single-command Docker Compose deployment
    \item RQ5 (Observability) — provided by Prometheus metrics and Grafana dashboards
\end{itemize}

The framework's modular design—separate processes communicating via HTTP, configurable parameters, and clear contract interfaces—makes it suitable for extension. Researchers can implement custom sequencers by following the interface specification in the repository documentation. A custom sequencer needs to: (1) expose HTTP endpoints for transaction submission and metrics reporting, (2) read contract addresses from the framework's configuration, and (3) call the contract's submit() function with the appropriate parameters.

The framework bridges the architectural understanding in Chapter~\ref{ch:framework} and the empirical evaluation in Chapter~\ref{ch:evaluation}. It operationalizes design objectives into concrete technical artifacts that can be deployed, tested, and analyzed.
